\documentclass[conference]{IEEEtran}

\usepackage[utf8]{inputenc}
\usepackage[T1]{fontenc}
\usepackage{newtxtext,newtxmath}
\usepackage{amsmath,amssymb,bm}
\usepackage{mathtools}
\usepackage{booktabs}
\usepackage{siunitx}
\usepackage{enumitem}
\usepackage[breaklinks=true,hidelinks]{hyperref}

\title{A Practical Description of the Schmitt-Trigger\\
Zero-Crossing Frequency Tracker and Its C\texttt{++} Implementation}

\author{%
  \IEEEauthorblockN{Mikhail S.\ Grushinskiy}
  \IEEEauthorblockA{Independent Researcher\\
  Bareboat Necessities Project\\
  \texttt{mgrouch@users.github.com}}
}

\begin{document}
\maketitle

\begin{abstract}
This note gives a clean and practical description of the zero-crossing
frequency tracker implemented in \texttt{SchmittTriggerZCFreqTracker}. The
method uses a Schmitt-trigger state machine, debounce timing, steepness
gating, and cycle-based period estimation. It is intended for lightweight
real-time tracking of approximately sinusoidal signals.

In addition to frequency, the implementation provides a phase-like estimate,
timing jitter, an amplitude-to-threshold ratio, a confidence score, and a
simple lock decision. The focus here is on what the code actually computes
and how the main parameters affect behavior.
\end{abstract}

\begin{IEEEkeywords}
frequency estimation, zero-crossing detector, Schmitt trigger, hysteresis,
debouncing, embedded signal processing
\end{IEEEkeywords}

% ---------------------------------------------------------------
\section{Introduction}

The tracker estimates the dominant frequency of a scalar oscillatory signal by
watching for alternating threshold crossings of a normalized input. Compared
with spectral methods or adaptive resonators, this approach is very simple:
it only needs a few state variables, a small period history buffer, and basic
timing logic.

The main design goals are:
\begin{itemize}[leftmargin=1.4em]
  \item reject chatter around zero using hysteresis,
  \item reject implausibly fast crossings using steepness and debounce timing,
  \item reduce jitter by averaging over multiple half-cycles,
  \item expose lightweight quality metrics suitable for embedded use.
\end{itemize}

% ---------------------------------------------------------------
\section{Inputs and Normalization}

At each update, the tracker receives:
\begin{itemize}[leftmargin=1.4em]
  \item the current signal sample \(s\),
  \item a positive signal-magnitude estimate \(M\),
  \item debounce time \(t_{\mathrm{deb}}\),
  \item steepness time \(t_{\mathrm{steep}}\),
  \item sample interval \(\Delta t\).
\end{itemize}

The implementation forms the normalized sample
\begin{equation}
  x = \frac{s}{|M|}.
\end{equation}

This is important: the tracker does \emph{not} estimate the signal magnitude
internally. It expects the caller to supply \(M\), and the hysteresis logic is
applied to the normalized quantity \(x\).

If \(\Delta t \le 0\) or \(|M| \le 0\), the update is skipped and the previous
frequency estimate is returned unchanged.

% ---------------------------------------------------------------
\section{Hysteresis Thresholds}

The tracker stores a positive hysteresis level \(h\) and uses the thresholds
\begin{equation}
  T_{\mathrm{up}} = +h,
  \qquad
  T_{\mathrm{low}} = -h.
\end{equation}

In the code,
\begin{equation}
  h = |\texttt{hysteresis}|,
\end{equation}
so the effective thresholds are always symmetric even if a negative value is
passed into the constructor.

The amplitude-to-threshold ratio reported by the tracker is
\begin{equation}
  R_A = \frac{|M|}{\max(h,10^{-12})}.
\end{equation}

This quantity does not affect crossing detection directly, but it influences
confidence and lock reporting.

% ---------------------------------------------------------------
\section{State Machine}

The tracker uses three states:
\begin{equation}
  \texttt{WAS\_NOT\_SET},\qquad
  \texttt{WAS\_LOW},\qquad
  \texttt{WAS\_HIGH}.
\end{equation}

\subsection{Startup State}

Initially the state is \texttt{WAS\_NOT\_SET}. In this state the tracker waits
for the normalized signal to move clearly outside the hysteresis band:
\begin{itemize}[leftmargin=1.4em]
  \item if \(x > T_{\mathrm{up}}\), the state becomes \texttt{WAS\_HIGH},
  \item if \(x < T_{\mathrm{low}}\), the state becomes \texttt{WAS\_LOW}.
\end{itemize}

At startup the internal timers and crossing counters are reset and the tracker
remains in fallback mode until valid crossings have been accumulated.

\subsection{Low and High States}

Once initialized, the tracker alternates between \texttt{WAS\_LOW} and
\texttt{WAS\_HIGH}. It keeps an internal running time
\begin{equation}
  t_{\mathrm{cyc}} = \texttt{timeInCycle},
\end{equation}
which is incremented by \(\Delta t\) at each update.

It also remembers the most recent time the signal was confidently on the low
or high side:
\begin{equation}
  t_{\mathrm{low,last}} = \texttt{lastLowTime},
  \qquad
  t_{\mathrm{high,last}} = \texttt{lastHighTime}.
\end{equation}

These stored times are used both for gating and for estimating crossing times.

% ---------------------------------------------------------------
\section{Crossing Detection}

\subsection{Low-to-High Crossing}

When the state is \texttt{WAS\_LOW}, a rising crossing is accepted only if:
\begin{equation}
  x > T_{\mathrm{up}},
\end{equation}
and
\begin{equation}
  t_{\mathrm{cyc}} - t_{\mathrm{low,last}} > t_{\mathrm{steep}},
\end{equation}
and either this is the first crossing in the current window or
\begin{equation}
  t_{\mathrm{cross}} - t_{\mathrm{cross,last}} > t_{\mathrm{deb}}.
\end{equation}

The crossing time estimate used by the code is
\begin{equation}
  t_{\mathrm{cross}}
  =
  t_{\mathrm{cyc}}
  - \frac{t_{\mathrm{cyc}} - t_{\mathrm{low,last}}}{2}.
\end{equation}

So the implementation uses a simple midpoint between the current sample time
and the last confirmed low-side time.

\subsection{High-to-Low Crossing}

Similarly, when the state is \texttt{WAS\_HIGH}, a falling crossing is
accepted only if:
\begin{equation}
  x < T_{\mathrm{low}},
\end{equation}
and
\begin{equation}
  t_{\mathrm{cyc}} - t_{\mathrm{high,last}} > t_{\mathrm{steep}},
\end{equation}
and either this is the first crossing in the window or
\begin{equation}
  t_{\mathrm{cross}} - t_{\mathrm{cross,last}} > t_{\mathrm{deb}}.
\end{equation}

The corresponding crossing-time estimate is
\begin{equation}
  t_{\mathrm{cross}}
  =
  t_{\mathrm{cyc}}
  - \frac{t_{\mathrm{cyc}} - t_{\mathrm{high,last}}}{2}.
\end{equation}

\subsection{Interpretation}

This is not a zero interpolation in the usual sense. The code does not fit a
line through neighboring samples around the true zero crossing. Instead, it
uses hysteresis and then places the crossing halfway between the current
sample time and the most recent time the signal was confidently on the
previous side of the Schmitt band.

% ---------------------------------------------------------------
\section{Period and Frequency Estimation}

\subsection{Crossing Window}

Let
\begin{equation}
  N_c = \texttt{crossingsCounter}
\end{equation}
be the number of accepted crossings in the current window. The first accepted
crossing sets the window start time
\begin{equation}
  t_{\mathrm{begin}} = \texttt{beginningCrossingInCycleTime}.
\end{equation}

Whenever a new accepted crossing occurs, the code forms
\begin{equation}
  T_{\mathrm{cycle}} = t_{\mathrm{cross}} - t_{\mathrm{begin}}.
\end{equation}

Since consecutive accepted crossings represent half-period events, the period
estimate used by the implementation is
\begin{equation}\label{eq:period_est}
  \hat T = \frac{2\,T_{\mathrm{cycle}}}{N_c - 1}.
\end{equation}

The frequency estimate is then
\begin{equation}
  \hat f = \frac{1}{\hat T}.
\end{equation}

\subsection{Configured Number of Periods}

The constructor parameter \texttt{periodsInCycle} controls how many full
periods are accumulated before a full update window is completed. If
\begin{equation}
  N_{\mathrm{per}} = \texttt{periodsInCycle},
\end{equation}
then a full window is reached when
\begin{equation}
  N_c = 2N_{\mathrm{per}} + 1.
\end{equation}

At that point the code computes \eqref{eq:period_est}, updates the period
statistics, updates the frequency, and restarts the local crossing window.

\subsection{Early Reacquisition Path}

The implementation also allows an earlier estimate before a full window is
completed. If more than one crossing has been observed and the current
frequency is still only the startup or fallback value, the code immediately
computes a provisional period and frequency using the same formula
\eqref{eq:period_est}. This helps the tracker recover faster after startup or
after long gaps with no valid crossings.

\subsection{Frequency Limits}

The code uses the constants
\begin{equation}
  f_{\mathrm{init}} = 10^{-3}\,\text{Hz},
\end{equation}
\begin{equation}
  f_{\mathrm{fallback}} = 10^{-2}\,\text{Hz},
\end{equation}
\begin{equation}
  f_{\max} = 10^{4}\,\text{Hz}.
\end{equation}

On full-window updates, the frequency is capped as
\begin{equation}
  \hat f \leftarrow \min(\hat f, f_{\max}).
\end{equation}

% ---------------------------------------------------------------
\section{Fallback Behavior}

If no accepted crossing occurs for too long, the tracker falls back to a low
default frequency. Let
\begin{equation}
  t_{\mathrm{fb}} = \texttt{fallbackToLowFreqTime},
\end{equation}
whose default is
\begin{equation}
  t_{\mathrm{fb}} = 60\,\text{s}.
\end{equation}

Whenever
\begin{equation}
  t_{\mathrm{cyc}} - t_{\mathrm{cross,last}} > t_{\mathrm{fb}},
\end{equation}
the code sets
\begin{equation}
  \hat f = f_{\mathrm{fallback}},
\end{equation}
and marks the estimate as fallback. This logic is active both after startup
and during normal operation if the signal stops producing valid alternating
crossings.

% ---------------------------------------------------------------
\section{Phase Estimate}

The tracker provides a phase-like diagnostic through
\texttt{getPhaseEstimate()}. If the frequency is non-positive or fallback mode
is active, it returns zero.

Otherwise, with period
\begin{equation}
  \hat T = \frac{1}{\hat f},
\end{equation}
and time since the last accepted crossing
\begin{equation}
  \Delta t_{\mathrm{cross}}
  =
  t_{\mathrm{cyc}} - t_{\mathrm{cross,last}},
\end{equation}
the reported phase is
\begin{equation}
  \phi
  =
  2\pi\,
  \operatorname{fmod}\!\left(
    \frac{\Delta t_{\mathrm{cross}}}{\hat T},\,1
  \right).
\end{equation}

If the current state is \texttt{WAS\_LOW}, the code adds a \(\pi\)-shift and
wraps back into \([0,2\pi)\):
\begin{equation}
  \phi \leftarrow (\phi+\pi)\bmod 2\pi.
\end{equation}

This is a useful phase-like timing indicator, but it should be viewed as a
diagnostic tied to the crossing logic rather than as a full analytic phase
estimate.

% ---------------------------------------------------------------
\section{Quality Metrics}

The tracker stores recent period estimates in a circular buffer of size
\begin{equation}
  K = 10.
\end{equation}

Let \(K_h\) be the number of valid entries currently stored, with
\(K_h \le K\). The code computes the population variance
\begin{equation}
  \sigma_T^2
  =
  \frac{1}{K_h}\sum_{k=1}^{K_h} T_k^2
  -
  \left(
    \frac{1}{K_h}\sum_{k=1}^{K_h} T_k
  \right)^2.
\end{equation}

Any small negative value caused by floating-point roundoff is clipped to zero.

The reported timing jitter is
\begin{equation}
  J = \sqrt{\max(\sigma_T^2,0)}.
\end{equation}

\subsection{Confidence}

If there is not yet enough period history, or if the most recent period is
non-positive, the confidence is set to
\begin{equation}
  c =
  \begin{cases}
    0, & \text{if fallback is active},\\
    0.1, & \text{otherwise}.
  \end{cases}
\end{equation}

Once enough valid period history exists, the code forms the normalized timing
spread
\begin{equation}
  r_T = \frac{J}{\max(T_{\mathrm{last}},10^{-12})},
\end{equation}
where \(T_{\mathrm{last}}\) is the most recent stored period estimate.

The base confidence is
\begin{equation}
  c_0 = \operatorname{clamp}_{[0,1]}(1-r_T).
\end{equation}

This is then shaped by the amplitude ratio:
\begin{equation}
  a_{\mathrm{score}}
  =
  \operatorname{clamp}_{[0,1]}\!\left(\frac{R_A}{2}\right),
\end{equation}
\begin{equation}
  c = c_0\,(0.35 + 0.65\,a_{\mathrm{score}}).
\end{equation}

If fallback mode is active, the final confidence is forced to zero.

\subsection{Reported Metrics}

The \texttt{QualityMetrics} structure returned by the code contains:
\begin{itemize}[leftmargin=1.4em]
  \item \texttt{confidence},
  \item \texttt{jitter},
  \item \texttt{amplitude\_ratio},
  \item \texttt{is\_fallback}.
\end{itemize}

\begin{table}[t]
  \centering
  \caption{Reported quality metrics}
  \label{tab:quality}
  \small
  \begin{tabular}{@{}lll@{}}
    \toprule
    Quantity & Code field & Meaning \\
    \midrule
    \(c\) & \texttt{confidence} & confidence score in \([0,1]\) \\
    \(J\) & \texttt{jitter} & period-jitter estimate in seconds \\
    \(R_A\) & \texttt{amplitude\_ratio} & magnitude relative to hysteresis \\
    fallback & \texttt{is\_fallback} & fallback-frequency flag \\
    \bottomrule
  \end{tabular}
\end{table}

% ---------------------------------------------------------------
\section{Lock Decision}

The tracker also exposes a boolean lock decision through \texttt{isLocked()}.
Lock is declared only when all of the following are true:
\begin{itemize}[leftmargin=1.4em]
  \item fallback mode is not active,
  \item at least two period estimates are available,
  \item confidence exceeds the configured threshold,
  \item amplitude ratio exceeds the configured minimum.
\end{itemize}

Equivalently,
\begin{equation}
  \texttt{locked}
  =
  \neg\texttt{isFallback}
  \land
  (\texttt{periodHistoryCount}\ge 2)
  \land
  (c \ge c_{\min})
  \land
  (R_A \ge R_{A,\min}).
\end{equation}

The defaults are
\begin{equation}
  c_{\min} = 0.60,
  \qquad
  R_{A,\min} = 1.0.
\end{equation}

These can be changed through \texttt{setLockParams()}.

% ---------------------------------------------------------------
\section{Implementation Notes}

The public interface of \texttt{SchmittTriggerZCFreqTracker} includes:
\begin{itemize}[leftmargin=1.4em]
  \item \texttt{update(...)} for per-sample processing,
  \item \texttt{getFrequencyHz()} for the latest frequency estimate,
  \item \texttt{getPhaseEstimate()} for the phase-like output,
  \item \texttt{getQualityMetrics()} for diagnostic metrics,
  \item \texttt{isLocked()} for the lock decision.
\end{itemize}

The implementation is intentionally simple. It does not require FFT buffers,
matrix operations, or long filters. Its quality depends strongly on the
supplied signal magnitude estimate and on reasonable choices of hysteresis,
debounce time, and steepness time.

% ---------------------------------------------------------------
\section{Summary}

The implemented Schmitt-trigger zero-crossing tracker is a compact real-time
frequency estimator with:
\begin{itemize}[leftmargin=1.4em]
  \item normalized hysteresis-based crossing detection,
  \item debounce and steepness gating for spurious-crossing rejection,
  \item half-cycle timing averaged into period estimates,
  \item fallback behavior for long crossing gaps,
  \item practical diagnostics for phase, jitter, confidence, and lock.
\end{itemize}

Its main strengths are simplicity, low computational cost, and robustness to
small fluctuations around zero. This makes it well suited to embedded
applications where a lightweight time-domain frequency tracker is preferred.

\end{document}
